% Drawing 144/08-2 -- folder 144, batch 08, archive page 153 (PDF page 154).
% Transcription: un grand cercle (une sphère) ; trois arcs descendent depuis le haut, issus d'un même voisinage du pôle, et s'arrêtent vers le milieu.
% Kind: sphere with arcs descending from near the pole. Method 2 (schematic redraw from the facsimile). Status: ok.
% Notes: A large circle; arcs come down from a neighbourhood of the top and stop near the middle. Disagreement (notes only): four such arcs are visible on the facsimile (the fourth close to the right edge of the circle), the transcription says three. A short stroke on the upper left of the circle is drawn.
% Coordinates: 1 unit = 1 mm = 1/100 of the crop width (bbox in page fractions: [0.08, 0.36, 0.52, 0.68]).
% Redrawn by Claude (claude-opus-5-5), issue #73, 2026-11-07; not checked by a human; not compiled (no TeX in the session).
% Rights: reproduces part of the work -- behind RENDER_TRANSCRIPTS, never published by itself.
\documentclass[tikz,border=3pt]{standalone}
\usepackage{amsmath,amssymb}
\usetikzlibrary{patterns,arrows.meta}
\tikzset{pencil/.style={black!45}, struck/.style={black!35}, mark/.style={circle,draw,inner sep=1.3pt}, pt/.style={circle,fill,inner sep=0.9pt}, opt/.style={circle,draw,fill=white,inner sep=1.1pt}}
\begin{document}
\begin{tikzpicture}[x=1mm,y=1mm,line cap=round,line join=round,font=\scriptsize]
\draw (55.9,-49.5) circle[radius=39.1];
\draw plot[smooth, tension=0.55] coordinates {(57.7,-12.2) (49.5,-28.6) (45.9,-46.8)};
\draw plot[smooth, tension=0.55] coordinates {(65.9,-13.2) (58.6,-28.6) (55.9,-47.7)};
\draw plot[smooth, tension=0.55] coordinates {(65.9,-14.1) (68.6,-28.6) (70.5,-43.2)};
\draw plot[smooth, tension=0.55] coordinates {(72.3,-12.2) (77.7,-24.1) (82.3,-42.2)};
\draw (25,-25.9) -- (28.6,-23.2);
\end{tikzpicture}
\end{document}
